\section{Experimental setup}\label{sec:setup}

\subsection{Datasets and client splits}\label{sec:setup-data}

\paragraph{Long-tailed benchmarks.} \R{main/datasets_cite} are made long-tailed as in GeFL-F's
code~\cite{kang2025gefl}: GeFL-F's fraction of the training set is divided over the classes by the profile of
\cref{sec:setting} with $\rho=100$, and no class receives more images than remain after the validation images are
held out. On CIFAR-10-LT the profile asks for more than that for the largest classes, which are therefore capped, so
the ratio between the largest and the smallest class is below 100 there (\cref{tab:datasets}); we name every setting
by its $\rho$. CIFAR-100-LT has 480 training images in its largest class and 5 in its smallest. The training images
are split among $K=10$ clients by a Dirichlet draw with $\alpha=0.5$ per class.
\paragraph{Natural imbalance.} \R{nat/datasets_cite} keep their own class sizes, with training sets capped at 20,000
images, and are split the same way. Training images that no client holds (200 per class, 20 per class on CIFAR-100,
and at most a quarter of a class with natural imbalance) form a validation set. It trains the referee classifier of the
generator diagnostics and is never used to select a model or a hyperparameter: every result is that of the last round,
on the official test set. \Cref{tab:datasets} summarises the splits.

\begin{table}[t]
\centering\small
\caption{Datasets as used: classes, training images held by the clients, test images, and the ratio between the
largest and the smallest class in the federation (natural imbalance: the dataset's own class sizes).}
\label{tab:datasets}
\slottab{datasets}{\begin{tabular}{lcccc}\toprule
Dataset & Classes & Training images & Test images & Largest / smallest class\\\midrule
\phrow{MNIST-LT}{4}\phrow{Fashion-MNIST-LT}{4}\phrow{SVHN-LT}{4}\phrow{CIFAR-10-LT}{4}\phrow{CIFAR-100-LT}{4}
\midrule
\phrow{DermaMNIST}{4}\phrow{OrganAMNIST}{4}\phrow{BloodMNIST}{4}\phrow{PathMNIST}{4}\phrow{GTSRB}{4}
\phrow{EMNIST (ByClass)}{4}\phrow{FER-2013}{4}\bottomrule\end{tabular}}
\end{table}

\subsection{Two protocols}\label{sec:setup-protocols}

\paragraph{Main setting: different client models.} As in GeFL-F, the clients share a one-block convolutional feature
extractor and each client has one of ten CNN head architectures of increasing depth and width (CNN-1 to CNN-10;
$K=10$, so one client per architecture). The rounds of the three stages are those of GeFL-F (e.g.\ 50, 200 and 100 on
CIFAR-10). All methods of this setting share the same Stage-I extractor for a given dataset and seed.

\paragraph{One shared model.} We follow the FedLTLib protocol~\cite{lin2026fedltlib}: CIFAR-10-LT and CIFAR-100-LT with
the standard long-tailed class profile $n_c=\lfloor n_{\max}\rho^{-c/(C-1)}\rfloor$ at $\rho\in\{100,50,10\}$, with
$n_{\max}$ the images per class left after 200 per class (20 on CIFAR-100) are held out for validation (rounding down
gives the smallest CIFAR-100 class 4 images at $\rho=100$), 20 clients, Dirichlet $\alpha=0.5$,
8 clients (40\%) per round, one ResNet-8~\cite{he2016resnet}, 200 rounds of 10 local epochs with batch size 32 and SGD
(learning rate 0.1, momentum 0.9, decay 0.99 per round). The baselines run in the authors' LongTailLib code, unchanged,
on our splits. For \ours the ResNet-8 is split after its second stage into the frozen extractor and the head; the
extractor is trained for 100 rounds and the head for 100 with 10 real epochs per round (the protocol's local epochs),
so that \ours uses the baselines' 200 rounds plus those of its generator, which we report separately. \ours averages
the clients' models with equal weights, as in the main setting, whereas the benchmark's FedAvg weights them by their
numbers of samples.

\subsection{Methods}\label{sec:setup-methods}

\Cref{tab:baselines} lists every method, its reference and the origin of its code. Methods designed for one shared
model are compared in the one-shared-model setting only. Local, in which each client trains alone, is a lower bound.
Hyperparameters are those of each method's paper; where the paper and the authors' code disagree we use the code.

\begin{table}[t]
\centering\small
\caption{Methods compared, their protocol and the origin of their code. HtFLlib~\cite{zhang2025htfllib} (Apache-2.0)
methods are ported to Flower with their authors' hyperparameters, except that they train for the 5 local epochs of
our head stage; LongTailLib~\cite{lin2026fedltlib} methods run in
the authors' code, downloaded at a fixed commit.}
\label{tab:baselines}
\begin{adjustbox}{max width=\linewidth}
\begin{tabular}{llll}
\toprule
Method & Venue & Protocol & Code\\
\midrule
GeFL-F~\cite{kang2025gefl} & IEEE TMC 2025 & different models & re-implementation of the authors' algorithm\\
Balanced Softmax~\cite{ren2020bsm} & NeurIPS 2020 & different models & loss on each client's head\\
FedLC~\cite{zhang2022fedlc} & ICML 2022 & both & loss re-implemented / LongTailLib\\
FedRS~\cite{li2021fedrs} & KDD 2021 & different models & loss re-implemented\\
FedDF~\cite{lin2020feddf} & NeurIPS 2020 & different models & re-implementation (distillation on generated features)\\
FedProto~\cite{tan2022fedproto} & AAAI 2022 & different models & HtFLlib\\
FedTGP~\cite{zhang2024fedtgp} & AAAI 2024 & different models & HtFLlib\\
FedGen~\cite{zhu2021fedgen} & ICML 2021 & different models & HtFLlib\\
FedGH~\cite{yi2023fedgh} & ACM MM 2023 & different models & HtFLlib\\
LG-FedAvg~\cite{liang2020lgfedavg} & arXiv 2020 & different models & HtFLlib\\
FedMRL~\cite{yi2024fedmrl} & NeurIPS 2024 & different models & HtFLlib\\
FedAvg~\cite{mcmahan2017fedavg}, FedProx~\cite{li2020fedprox}, MOON~\cite{li2021moon} & & one shared model & LongTailLib\\
CCVR~\cite{luo2021ccvr}, CReFF~\cite{shang2022creff}, FedETF~\cite{li2023fedetf} & & one shared model & LongTailLib\\
Fed-GraB~\cite{xiao2023fedgrab}, FedLoGe~\cite{xiao2024fedloge} & & one shared model & LongTailLib\\
FedYoYo~\cite{yan2025fedyoyo}, RUCR~\cite{huang2024rucr} & & one shared model & LongTailLib\\
Local & -- & different models & each client alone\\
\ours & this work & both & ours (Flower)\\
\bottomrule
\end{tabular}
\end{adjustbox}
\end{table}

\subsection{Metrics and statistics}\label{sec:setup-metrics}

The primary metric is balanced accuracy on the test set, with the heads as trained and with the exact readout. We
also report macro-F1, macro precision, the multi-class Matthews correlation coefficient (MCC)~\cite{gorodkin2004mcc},
the geometric mean of the class recalls (each floored at 1\%), the recall of the tail classes (the 40\% least
frequent) and of the worst class, macro AUC-PR, the expected calibration error with 15 bins
(ECE)~\cite{guo2017calibration} and its class-wise form~\cite{kull2019classwise}, and, for the main setting, the balanced
accuracy of the weakest architecture and the accuracy on each client's own class mix. In the setting with different
client models every metric is, as in GeFL~\cite{kang2025gefl}, the mean over the ten client models (one per
architecture), each evaluated on the whole test set; the geometric mean and the worst-class recall are computed for
each model and then averaged, which makes them lower than the same metrics on the models' pooled predictions. Every
per-class metric is
logged for every class at every evaluation. Cost is the measured communication per client, the wall-clock time on one
GPU and the floating-point operations.

Every configuration is run with seeds 3, 4 and 5; a seed fixes the split and the initialisations of every method, so
runs with the same seed are paired. We report mean $\pm$ sample standard deviation over seeds; for differences to
\ours, the mean paired difference with a 95\% $t$ interval, the paired $t$-test and Holm's correction~\cite{holm1979}
over the compared methods. Settings that were run with one seed say so and carry no interval.

\subsection{Implementation}\label{sec:setup-impl}

\ours and every method of the main setting run in Flower 1.42~\cite{beutel2020flower} (simulation with the Ray
backend), one simulation per run; the one-shared-model baselines run in the authors' LongTailLib code on the same
client splits and are evaluated on the same test set; all runs are on
NVIDIA T4 GPUs. All releases of \ours use Flower's SecAgg+~\cite{bell2020secaggplus} with 30-bit quantisation, every
vector divided by a public bound before quantisation. The Flower implementation reproduces the client splits of the
original GeFL-F implementation exactly and its accuracy within seed noise (\cref{app:protocol}). \Cref{tab:hyper}
lists the settings of \ours, which are the same on every dataset.

\begin{table}[t]
\centering\small
\caption{Settings of \ours (all datasets; with one shared model, 10 real epochs per round, the protocol's local
epochs).}
\label{tab:hyper}
\begin{adjustbox}{max width=\linewidth}
\begin{tabular}{ll}
\toprule
Setting & Value\\
\midrule
Separation gate threshold $F_0$ & 0.05\\
CCWA parameter $\beta$ & 0.999\\
Moment correction: samples per class & 500\\
Moment correction: range of $s_c$ & $[0.1, 10]$\\
Synthetic / real epochs per round $T_s$ / $T_r$ & 10 / 5\\
Logit adjustment $\tau$ & 1\\
Consensus soft-label weight $\beta$ & 0.5\\
Readout shrinkage $\gamma$ & 0.1\\
Generator optimiser, learning rate & Adam, $10^{-3}$\\
Generator weight decay (coupled L2) & $10^{-3}$\\
\bottomrule
\end{tabular}
\end{adjustbox}
\end{table}
